\documentclass[../../../include/open-logic-section]{subfiles}

\begin{document}

\olfileid{sth}{z}{union}
\olsection{યોગગણ}

\olref[sth][z][sep]{prop:intersectionsexist}માંથી આપણને છેદગણો
મળ્યા. પરંતુ મનસ્વી યોગગણો અસ્તિત્વ ધરાવે એવું ઇચ્છીએ, તો
બીજી સ્વયંસિદ્ધિ આપવી પડે:

\begin{axiom}[યોગગણ]
કોઈ પણ ગણ $A$ માટે ગણ $\bigcup A =
\Setabs{x}{(\exists b \in A) x \in b}$ અસ્તિત્વ ધરાવે છે.
\[
\forall A \exists U \forall x(x \in U \liff (\exists b \in A)x \in b)
\]
\end{axiom}

આ સ્વયંસિદ્ધિ પણ સંચયી-પુનરાવર્તી કલ્પનાથી યોગ્ય ઠરે છે.
$A$ને ગણ લો; તેથી $A$ કોઈ તબક્કા $S$ પર રચાય છે
(\stageshier પ્રમાણે). $A$નો દરેક સભ્ય $S$ની \emph{પહેલાં}
રચાયો હતો (\stagesacc પ્રમાણે); તેથી એ જ રીતે વિચારતાં
$A$ના દરેક સભ્યનો દરેક સભ્ય $S$ પહેલાં રચાયો હતો.
આથી \emph{તે બધા} ગણો $S$ પહેલાં ઉપલબ્ધ છે, જેથી
$S$ પર તેમનો ગણ રચી શકાય. અને તે ગણ બરાબર $\bigcup A$ છે.

\end{document}
